% paper/sections/05_ccd.tex
\section{結合コンパクト差分法（CCD 法）の厳密導出}
\label{sec:CCD}

\subsection{CCD スキームの定義}

CCD 法は，関数値 $f$ とその一次微分 $f^{(1)}$，二次微分 $f^{(2)}$ を同時に連立して解くことで，
3点ステンシルで $\Ord{h^6}$ 精度を実現する手法である．
各内点 $i$ において，以下の2式が成立するように係数を定める：

\begin{tcolorbox}[mybox, title={CCD スキームの2方程式}]
\textbf{Equation-I}（一次微分の結合式）：
\begin{equation}
  \alpha_1\fp{i-1} + \fp{i} + \alpha_1\fp{i+1}
  = \frac{a_1}{h}\bigl(f_{i+1}-f_{i-1}\bigr)
  + b_1 h\bigl(\fpp{i+1}-\fpp{i-1}\bigr)
  \label{eq:CCD_I}
\end{equation}

\textbf{Equation-II}（二次微分の結合式）：
\begin{equation}
  \beta_2\fpp{i-1} + \fpp{i} + \beta_2\fpp{i+1}
  = \frac{a_2}{h^2}\bigl(f_{i-1}-2f_i+f_{i+1}\bigr)
  + \frac{b_2}{h}\bigl(\fp{i+1}-\fp{i-1}\bigr)
  \label{eq:CCD_II}
\end{equation}
\end{tcolorbox}

\subsection{Equation-I の係数導出}

各項を格子点 $i$ の周りでテイラー展開し，次数ごとに整理することで係数を決定する．
例えば $f_{i\pm1} = f_i \pm h f^{(1)}_i + \frac{h^2}{2} f^{(2)}_i \pm \frac{h^3}{6} f^{(3)}_i + \dots$ を用いる．

Equation-I の左辺 (LHS) と右辺 (RHS) を展開すると以下のようになる：
\begin{align}
  \text{LHS} &= (1+2\alpha_1)\fp{i} + \alpha_1 h^2 f^{(3)}_i + \frac{\alpha_1}{12}h^4 f^{(5)}_i + \Ord{h^6} \\
  \text{RHS} &= \left(2a_1 \fp{i} + \frac{a_1}{3}h^2 f^{(3)}_i + \frac{a_1}{60}h^4 f^{(5)}_i\right) \notag \\
             &\quad + \left(2b_1 h^2 f^{(3)}_i + \frac{b_1}{3}h^4 f^{(5)}_i\right) + \Ord{h^6}
\end{align}

両辺の各次数の係数を比較し，誤差項を $O(h^4)$ まで消去する連立方程式を立てる：
\begin{itemize}
  \item $\fp{i}$ の係数： $1 + 2\alpha_1 = 2a_1$
  \item $f^{(3)}_i$ の係数 ($h^2$)： $\alpha_1 = \frac{a_1}{3} + 2b_1$
  \item $f^{(5)}_i$ の係数 ($h^4$)： $\frac{\alpha_1}{12} = \frac{a_1}{60} + \frac{b_1}{3}$
\end{itemize}
この3元連立1次方程式を解くことで，$\Ord{h^6}$ 精度を達成する係数が一意に定まる．

\begin{tcolorbox}[resultbox]
\textbf{Equation-I 係数解：}
\begin{equation}
  \boxed{
    \alpha_1 = \frac{7}{16},\quad
    a_1 = \frac{15}{16},\quad
    b_1 = \frac{1}{16}
  }
\end{equation}
\end{tcolorbox}

\subsection{Equation-II の係数導出}

同様に Equation-II についても，$f^{(2)}_i$ を中心に展開を行う．
\begin{itemize}
  \item $f^{(2)}_i$ の係数： $1 + 2\beta_2 = a_2 + 2b_2$
  \item $f^{(4)}_i$ の係数 ($h^2$)： $\beta_2 = \frac{a_2}{12} + \frac{b_2}{3}$
  \item $f^{(6)}_i$ の係数 ($h^4$)： $\frac{\beta_2}{12} = \frac{a_2}{360} + \frac{b_2}{60}$
\end{itemize}
これらを解くことで，以下の係数を得る．

\begin{tcolorbox}[resultbox]
\textbf{Equation-II 係数解：}
\begin{equation}
  \qquad
  \boxed{
    \beta_2 = -\frac{1}{8},\quad
    a_2 = 3,\quad
    b_2 = -\frac{9}{8}
  }
  \label{eq:coef_CCD}
\end{equation}
\end{tcolorbox}

\subsection{境界条件}
\label{sec:ccd_bc}

CCD スキームは内点において3点ステンシル（$i-1, i, i+1$）を用いるため，両端の境界点（$i=0, N$）では方程式が閉じない問題が生じる．
この「閉包問題」を解決するため，境界点専用の片側差分スキームが必要となる．

\begin{tcolorbox}[warnbox, title={境界スキームの役割と精度}]
境界スキームは，内点の連立方程式系（ブロック三重対角行列）を閉じるために不可欠である．
本稿では，数値的な安定性と精度のバランスを考慮し，内点より1次低い $\Ord{h^5}$ 精度のスキームを採用する．
具体的には，境界点 $i=0$ での微分値 $\fp{0}, \fpp{0}$ を，境界近傍の関数値 $f_0, \dots, f_3$ と，内側の微分値 $\fp{1}, \fpp{1}$ を用いて表現する式を導出する．
これにより，全体としての計算精度を損なうことなく，ロバストな境界処理が可能となる．
\end{tcolorbox}

左境界 ($i=0$) における Equation-I 相当の結合式は以下のようになる：
\begin{equation}
  \fp{0} + \frac{3}{2}\fp{1} - \frac{3h}{2}\fpp{1}
  = \frac{1}{h}\!\left(-\frac{23}{6}f_0 + \frac{21}{4}f_1 - \frac{3}{2}f_2 + \frac{1}{12}f_3\right)
\end{equation}

\subsection{非一様格子への拡張：座標変換法}
\label{sec:ccd_metric}

本稿では界面近傍の解像度を上げるために非一様格子を用いるが，前述の CCD 係数（式 \eqref{eq:coef_CCD}）は等間隔格子 $h$ を前提としている．
非一様格子上で高精度を維持するための最もエレガントな手法は，計算空間への座標変換を用いることである．

\begin{tcolorbox}[mybox, title={座標変換による実装}]
\textbf{概念：}
物理空間 $x$（非一様）を，計算空間 $\xi$（一様，$\Delta\xi=1$）に写像する．
微分計算はすべて計算空間 $\xi$ 上で標準の等間隔 CCD スキームを用いて行い，その結果を連鎖律（chain rule）で物理空間の微分値に変換する．

\textbf{変換公式：}
物理空間での一次微分 $\partial f/\partial x$ および二次微分 $\partial^2 f/\partial x^2$ は，メトリック係数 $J = \partial \xi / \partial x$ を用いて次のように計算される：
\begin{align}
  \pd{f}{x} &= \pd{\xi}{x}\pd{f}{\xi} = J \fp{\xi} \\
  \pdd{f}{x} &= J \pd{}{x}\!\left(\fp{\xi}\right) + \fp{\xi}\pd{J}{x} = J^2 \fpp{\xi} + J \pd{J}{\xi} \fp{\xi}
\end{align}
ここで，$f^{(\cdot)}_\xi$ は計算空間上で CCD 法により求めた微分値である．
この手法により，複雑な非一様格子係数を導出することなく，一様格子の CCD ソルバーをそのまま再利用できる利点がある．
\end{tcolorbox}
